%! Introducing Statistics: What Can We Learn from Data? — Unit 1, Lesson 1.0 — Group Activity (Answer Key)
\documentclass[10pt]{article}
\usepackage{apstats-article}
\usepackage{apstats-key}   % pulls in -boxes; do NOT also load -boxes

\begin{document}

\pageheader{Unit 1, Lesson 1.0}{Group Activity --- Answer Key}

\vspace{0.15in}

\boxguard[14]
\begin{headlinebox}{goldbg}
\small
\textbf{Your group works one tier.} Every tier ends the same way: with a statistical
question you wrote yourselves. Be ready to read it to the class and defend it with the
two checks --- \emph{do the answers differ?} and \emph{must we collect data?}
\end{headlinebox}

\vspace{0.14in}

\boxguard[30]
\begin{scenariobox}[Tier R --- Read the Data Set]{navy}
\small
A study hall keeps this record for one week:

\vspace{6pt}
\renewcommand{\arraystretch}{1.35}
\begin{tabularx}{\linewidth}{@{} l Y Y Y @{}}
\rowcolor{sky}
\textbf{Student} & \textbf{Grade} & \textbf{Homework (min)} & \textbf{Has a job?} \\
\hline
Ava    & 10 & 45 & No  \\
Ben    & 11 & 20 & Yes \\
Cruz   & 10 & 75 & No  \\
Dee    & 12 & 30 & Yes \\
Eli    & 11 & 60 & No  \\
\hline
\end{tabularx}

\vspace{10pt}
\textbf{1.} One individual in this data set is \ans{one student, e.g.\ Ava}.

\vspace{6pt}
\textbf{2.} There are \ans{5} individuals and \ans{4} variables.

\vspace{6pt}
\textbf{3.} ``11'' is a \ans{value} of the variable \ans{Grade}.

\vspace{8pt}
\textbf{4.} \textbf{Circle} the statistical question.
\begin{enumerate}[label=\textbf{(\Alph*)}, leftmargin=2em, itemsep=4pt]
  \item How many minutes of homework did Cruz do?
  \item Is Dee in the twelfth grade?
  \item How do homework times vary among these students?
        \textcolor{keyred}{\textbf{$\leftarrow$ correct}}
\end{enumerate}

\vspace{6pt}
\textbf{5.} Why do the other two fail the test?\\[4pt]
\ansline{Each asks about a single individual, so each has one definite answer ---}\\[1pt]\ansline{nothing varies, and the value can be read off the table without collecting data.}
\end{scenariobox}

\vspace{0.14in}

\boxguard[26]
\begin{scenariobox}[Tier A --- Write the Questions]{greenacc}
\small
At football tryouts, a coach times every player in the 40-yard dash and records the
time in seconds, along with the position each player is trying out for.

\vspace{10pt}
\textbf{1.} The individuals are \ans{the players at tryouts}.

\vspace{6pt}
\textbf{2.} Name the two variables: \ans{40-yard dash time} and \ans{position}.

\vspace{8pt}
\textbf{3.} Write \textbf{two} statistical questions the coach could answer with this
data.\\[4pt]
\ansline{``How do the 40-yard dash times of the players vary?'' and ``Do dash times}\\[1pt]\ansline{differ across the positions?'' Both anticipate answers that differ by player.}

\vspace{8pt}
\textbf{4.} For your first question, what exactly \emph{varies} from one individual to
the next? Be specific.\\[4pt]
\ansline{The dash time, in seconds. Each player runs his own time, and those times}\\[1pt]\ansline{are not all the same --- that spread is exactly what the question asks about.}

\vspace{8pt}
\textbf{5.} Write one question about these players that is \emph{not} statistical, and
say what makes it fail.\\[4pt]
\ansline{``What was Player 12's dash time?'' It names one individual, so it has one}\\[1pt]\ansline{definite answer --- nothing varies and no new data must be collected.}
\end{scenariobox}

\vspace{0.14in}

\boxguard[26]
\begin{scenariobox}[Tier E --- Interrogate the Headline]{goldacc}
\small
A news site runs this headline and nothing else:

\vspace{6pt}
\begin{center}
{\large\bfseries\color{navy} ``Teens spend 7 hours a day on screens.''}
\end{center}

\vspace{8pt}
\textbf{1.} List \textbf{three} pieces of context the headline leaves out.\\[4pt]
\ansline{1. WHO --- which teens, from where, and how many of them were measured?}\\[1pt]\ansline{2. WHAT --- which screens count? Phones only, or TV, laptops, and games too?}\\[1pt]\ansline{3. WHEN and HOW --- what year, and were the hours self-reported or tracked?}

\vspace{8pt}
\textbf{2.} The number 7 is a single value, but it summarizes many individuals. What
was the variable, and who were the individuals?\\[4pt]
\ansline{Individuals: the teens in the study. Variable: hours of screen time per day,}\\[1pt]\ansline{recorded for each teen. The 7 is one summary of all those differing values.}

\vspace{8pt}
\textbf{3.} Write the statistical question the study behind this headline must have
asked.\\[4pt]
\ansline{``How much daily screen time do teens get, and how does it vary among them?''}\\[1pt]\ansline{The answers differ from teen to teen, so data had to be collected to find them.}

\vspace{8pt}
\textbf{4.} A classmate says, ``So every teen is on a screen 7 hours a day.'' Explain
in one sentence why that misreads the number.\\[4pt]
\ansline{The 7 summarizes the whole group; individual teens are above and below it,}\\[1pt]\ansline{and the summary says nothing about any one teen's screen time.}
\end{scenariobox}

\end{document}
